\subsection{A bounded negative moment outside the finite-level de Rham condition}
\label{sec:cubic-inverse-period-limit}

This is an explicit local counterexample to a proposed general limit
inference.  It does not compute the actual Rankin--Iwasawa classes of
AC-K33-image and does not assert their vanishing or nonvanishing.

\begin{proposition}[AC-K34: finite-level de Rham conditions do not pass through every bounded negative moment]
\label{AC-K34}
Let \(p>2\), \(K=\mathbb Q_p(\mu_p)\), and
\(K_\infty=K(\mu_{p^\infty})\).  Put
\(\Gamma=\operatorname{Gal}(K_\infty/K)\).
Write \(\chi\) for the cyclotomic character, put \(u=1+p\),
and choose the topological generator \(\gamma\) with
\(\chi(\gamma)=u\).  Then
\[
 a(g)=\frac{\log\chi(g)}{\log u}\in\mathbb Z_p,\qquad
 \overline g=\gamma^{a(g)}.
\]
Set \(\Lambda=\mathbb Z_p[[\Gamma]]=\mathbb Z_p[[X]]\),
where \(X=\gamma-1\).  Give \(\Lambda\) the left regular action
\(g(m)=\gamma^{a(g)}m\), and give \(\Lambda(1)\) the action
\[
 g(m)=u^{a(g)}\gamma^{a(g)}m .
\]
All topologies are the compact inverse-limit topology, equivalently
the \((p,X)\)-adic topology on \(\Lambda\).
There are an integral continuous cocycle \(z\) with coefficient
\(\Lambda(1)\) and an integral Galois-equivariant map
\(q:\Lambda(1)\to\mathbb Z_p\) such that every finite
permutation-level image of \(z\) is a rational coboundary, whereas
\(q_*[z]\notin H_g^1(K,\mathbb Q_p)\).
The map \(q\) is bounded for the integral measure norm and continuous
for the compact \((p,X)\)-adic topology.
Its unique finite-level rational primitives are compatible but require
unbounded denominators, of exact valuation \(n+1\) at level \(p^n\).
\end{proposition}

\begin{proof}
\proofstep{An explicit integral cocycle and an integral negative moment.}
The formula
\[
 c(g)=\frac{\gamma^{a(g)}-1}{\gamma-1}
      =\sum_{k\geq0}\binom{a(g)}{k+1}X^k
\]
defines a continuous integral cocycle with coefficient \(\Lambda\).
Indeed, each binomial coefficient is continuous in \(a(g)\), and
\[
 c(gh)=c(g)+\gamma^{a(g)}c(h).
\]
Its augmentation is \(a(g)\).

Define the continuous \(\mathbb Z_p\)-algebra automorphism
\[
 A:\Lambda(1)\longrightarrow\Lambda,\qquad
 A(\gamma)=u^{-1}\gamma .
\]
Here \(A\) denotes an automorphism of the underlying compact module
with the indicated, different Galois actions.  It is well-defined
because \(u^{-1}(1+X)-1\in(p,X)\); its inverse sends
\(\gamma\) to \(u\gamma\).
It is Galois-equivariant, since
\[
 A\bigl(u^{a(g)}\gamma^{a(g)}m\bigr)
   =\gamma^{a(g)}A(m).
\]
Consequently \(z=A^{-1}c\) is a continuous integral cocycle with
coefficient \(\Lambda(1)\).  Explicitly,
\[
 z(g)=\frac{(u\gamma)^{a(g)}-1}{u\gamma-1}.
\]
This expression means its convergent binomial series at
\(u\gamma-1\), not an inversion in the integral Iwasawa algebra.

The map
\[
 q=\operatorname{aug}\circ A:\Lambda(1)\longrightarrow\mathbb Z_p
\]
is a bounded, Galois-equivariant negative-character moment.
In measure notation it integrates \(\chi^{-1}\); in power-series
notation it evaluates \(\gamma\) at \(u^{-1}\).
For the measure norm
\(\|\mu\|=\sup_C|\mu(C)|_p\), with \(C\) ranging over clopen
subsets of \(\Gamma\), multiplication by the continuous unit-valued
character \(\chi^{-1}\) and its inverse preserve the integral unit ball.
Thus \(A\) is an isometry for this norm, and \(q\) has norm at most one.
This norm assertion is separate from compact-topology continuity,
which follows from the displayed power-series substitutions.
It satisfies
\[
                         q(z(g))=a(g).
\]

\proofstep{Every finite-level rational class vanishes.}
Let
\[
 \Lambda_n=\mathbb Z_p[\Gamma/\Gamma^{p^n}],\qquad
 Q_n=p^n,
\]
with the induced twisted action on \(\Lambda_n(1)\).
In \(\mathbb Q_p\otimes\Lambda_n\), put
\[
 b_n=
 \frac{\displaystyle\sum_{r=0}^{Q_n-1}u^r\gamma^r}
      {u^{Q_n}-1}.
\]
The identity \(\gamma^{Q_n}=1\) gives
\[
 (u\gamma-1)b_n=1.
\]
Thus \(u\gamma-1\) is invertible on this finite rational module,
and
\[
 z_n(g)=(g-1)b_n .
\]
Every \(z_n\) is therefore a rational coboundary, and in particular
its class belongs to \(H_g^1(K,\mathbb Q_p\otimes\Lambda_n(1))\).
Alternatively finite Shapiro identifies this de Rham condition
with that on \(\mathbb Q_p(1)\) over
\(K_n=K(\mu_{p^{n+1}})\); Bloch--Kato Example 3.9 gives
\(H_g^1(K_n,\mathbb Q_p(1))=H^1(K_n,\mathbb Q_p(1))\).
The explicit coboundary proof is stronger and needs no inverse-limit
cohomology comparison.

The primitive \(b_n\) is unique: a Galois invariant would be killed
by the invertible operator \(u\gamma-1\).
The natural transition satisfies
\[
                 \pi_{n+1,n}(b_{n+1})=b_n .
\]
This follows either by the finite geometric-series formula or by
uniqueness, since both sides are primitives of \(z_n\).

For odd \(p\),
\[
                      v_p(u^{p^n}-1)=n+1.
\]
Every numerator coefficient in the displayed basis
\(1,\gamma,\ldots,\gamma^{p^n-1}\) is a unit.
Hence no fixed power of \(p\) clears all the \(b_n\).
The compatible family belongs to
\[
 \varprojlim_n(\mathbb Q_p\otimes\Lambda_n(1))
\]
but not to the image of
\[
          \bigl(\varprojlim_n\Lambda_n(1)\bigr)[1/p].
\]
This explicitly locates the failure of interchanging rationalization
with this inverse limit.

\proofstep{The bounded completed-tower image is not de Rham.}
The class \(q_*[z]\) is represented by the nonzero ramified homomorphism
\(a:G_K\to\mathbb Q_p\).
Bloch--Kato Example 3.9, printed p.~359, states that
\(H_g^1(K,\mathbb Q_p)=H_f^1(K,\mathbb Q_p)\) is the one-dimensional
space of unramified homomorphisms.
Since \(a\) is nonzero on cyclotomic inertia,
\[
                         q_*[z]\notin H_g^1(K,\mathbb Q_p).
\]

There is no contradiction to finite-dimensional functoriality.
The map \(q\) does not factor through any finite \(\Lambda_n\):
\[
 q(\gamma^{p^n}-1)=u^{-p^n}-1\ne0.
\]
It also cannot be extended to a Galois-equivariant map on the whole
inverse system of rational primitives while preserving the indicated
coefficient map, since such an extension would make \(q_*[z]\) a
coboundary.  A fortiori, finite-level period primitives alone do not
supply the period primitive required after the completed-tower map.
\end{proof}

\paragraph{Exact relevance to AC-K33-image.}
The compact finite-coefficient arguments in AC-K33-image correctly
construct the global classes and their bounded coefficient images.
They do not assert a uniform bound for rational or \(B_{\rm dR}\)
primitives, and none is supplied by the boundedness of the coefficient
map itself.  A finite-level representative modulo \(p^r\) of a
negative-moment vector is not an exact characteristic-zero
Galois-equivariant map at that finite level.

The example proves only that those facts, together with finite-level
geometricity or membership in \(H_g\), do not logically imply that
the actual K33 image lies in \(H_g\).  The specific Rankin classes may
satisfy additional identities not present here.  No singular-rank
claim follows in either direction.  An actual vanishing argument
would still need compatible period primitives in a specified
coefficient-period topology through which the actual quotient
extends; an actual nonvanishing argument would still need the
missing reciprocity calculation.

\paragraph{Primary-source scope.}
The only p-adic Hodge input in the counterexample is
Bloch--Kato, \emph{Tamagawa numbers of motives}, Definition (3.7.2)
and Example 3.9, printed pp.~353 and 359:
\url{https://virtualmath1.stanford.edu/~conrad/BSDseminar/refs/BKTamagawa.pdf}.
These passages and the statement of Proposition 3.8 were directly
reread for this note.  The complete Proposition 3.8 proof is retained
at the already-audited AC-G69 scope, not claimed as a new independent
proof here.  All cocycle, twist, denominator and inverse-limit
calculations above are supplied explicitly.

